\section{The Deployment Front-Door Rule}
\label{sec:rule}

The attribution gap calls for a limited information rule rather than a new general law of AI liability. The rule must tell an affected person whom to approach, tell enterprises what record to maintain, and tell courts when and how that record becomes available. It must also provide a prompt exit for an enterprise whose connection to the brand or technical stack did not include practical control over the asserted risk.

This Part defines sectoral coverage, the visible-operator gateway, material control, staged access, protection, remedies, and safe exit. It then supplies model statutory text and fixes the boundary between the ex ante state-law duty and federal party discovery.

\subsection{Sectoral Designation Defines the Field}
\label{sec:designation}

``High impact'' should identify a deployment category, not attach permanently to a model. A general-purpose model can support casual drafting in one setting and become part of a benefits, hiring, clinical, or physical-control process in another. The deployment supplies the protected interest, affected population, operating institution, and legally relevant control. Coverage should follow those features.

The proposed statute therefore applies only when a legislature has specified the protected interests, visible operator, minimum record, retention and material-change rules, claimant threshold, and remedy. It may designate a category directly or authorize a sectoral agency to maintain a bounded list under express statutory criteria; the legislature retains the choices that determine who must answer and which interests justify compulsory recordkeeping.

Five features identify the strongest candidates. First, the deployment affects an interest already protected by substantive law, such as employment opportunity, access to public benefits, health, credit, housing, education, physical safety, or the welfare of minors. Second, the system repeatedly mediates decisions or exposure rather than producing an isolated low-consequence convenience. Third, affected people predictably observe less about the operative system than the operator and suppliers know. Fourth, the identity, version, control, and event fields can be standardized at reasonable cost. Fifth, existing doctrine can adjudicate the dispute once those facts are available. These criteria select the settings in which a record gateway unlocks an established legal regime rather than inventing a new one.

Designation also fixes the public-law interface. If a covered operator is a government institution, the legislature should specify the enforcing tribunal, available relief, applicable exhaustion route, and the treatment of sovereign or official immunity. The front door need not imply a damages action against the government. It can supply an administrative record, declaratory remedy, reviewable denial, or other access path tailored to the underlying interest. Defining that route at designation keeps the information right aligned with the sector's remedial structure.

This design avoids an unstable universal definition. Employment screening, allocation of essential public benefits, clinical diagnosis or treatment, credit, housing, education admissions, physical safety, and relational services offered to minors may justify different records and retention periods. A companion service may require age-access, relational-behavior, crisis-intervention, and continuity records. An employment tool may require input-feature, validation, selection-rate, configuration, and version records. A clinical system may require intended-use, validation-population, uncertainty, escalation, override, and follow-up records. The shared front-door architecture permits sector-specific content.

California's enacted companion-chatbot provisions demonstrate the feasibility of this form. The legislature defined a covered relational service and its operator, required specified disclosures and self-harm protocols, imposed protections for known minors and reporting obligations, supplied a private remedy, and preserved cumulative law.\lrfootnote{See \statcite{californiaSB2432025}. The enacted provisions do not contain the deployment front-door rule proposed here; they demonstrate that a legislature can designate a relational deployment, identify its operator, and attach enforceable protocol and reporting duties.} The statute did not need to declare every conversational model high impact. It regulated a deployment category and the enterprise that made it available.

A designation should also identify exclusions. Research or internal evaluation that does not produce a legally or practically consequential effect on an affected person ordinarily falls outside. So does general cloud hosting without risk-specific configuration or intervention authority. Small entities may receive simplified record forms or regulator assistance where the sector permits, but the information essential to identify a consequential decision should not disappear with enterprise size. A covered operator using open weights still controls a deployment; an upstream author who has no continuing authority may fall outside the post-release control lane.

\subsection{The Visible Operator Answers First}
\label{sec:gateway}

The rule resolves the service problem by selecting the entity an affected person can ordinarily identify. The \term{visible operator} is the enterprise or public institution that made the covered deployment available to the affected person or incorporated it into the process that produced the adverse event. It may be a companion-service provider, employer, hospital, lender, benefits agency, school, or operator of a branded application. Its gateway role follows from visibility and provenance, not a presumption that it controlled every risk.

Stage One uses an occurrence threshold rather than a miniature merits pleading. A verified notice identifies: (1) an adverse event affecting an interest protected by the law specified in the sectoral designation; (2) the covered deployment or visible operator through which the event occurred; (3) facts reasonably connecting the requested identity and event fields to that occurrence; (4) the relevant time, transaction, account, decision, or interaction with as much specificity as reasonably available; (5) the closed fields sought; and (6) a good-faith intended use of the answer to evaluate, contest, appeal, or seek relief for the event. A declaration under penalty of perjury supplies the requester's personal knowledge and marks matters stated on information reasonably available; counsel certifies proper purpose and reasonable inquiry.

That threshold demands an event and a bounded connection, but never a fact held exclusively in the deployment record. An applicant need not know whether a named feature was enabled, whether it generated a score, whether the score reached the employer, which version ran, or which supplier retained the event log. Those are first-answer fields. ``An AI harmed me'' remains insufficient because it identifies neither a covered occurrence nor a closed record. But the operator cannot defeat Stage One by demanding proof that an undisclosed feature participated in the event.

The notice invokes a statutory access right that exists outside merits discovery and may be served before a complaint. A designation may also permit delivery in another lawful forum; model Act section 7 supplies the separate boundary after a federal action concerning the event is pending. Delivery follows the channel named in the designation, invokes the first-answer duty, and is not service of process for a later action. The operator receives thirty days, subject to one short documented extension, to answer, object by field, or state that no covered deployment processed the event. Silence, a blanket proprietary designation, and referral to an unnamed vendor are not responses.

If the operator refuses or materially underanswers, the affected person may seek narrow enforcement against that known entity in a lawful forum. The statutory violation is complete upon coverage, a compliant request, delivery, and nonperformance; enforcement does not depend on pleading the merits fact that the first answer exists to reveal. The administering agency or a court with independent subject-matter and personal jurisdiction can order the defined answer, and a valid arbitration agreement may select the private enforcement forum. The requester bears the burden on coverage, delivery, and the occurrence threshold. The operator bears the burden on a field-specific objection based on facts within its control, including nonoperation, burden, privilege, privacy, security, or documented unavailability. Jurisdiction, immunity, preemption, and a pure legal bar remain available before access expands beyond the fixed first answer.

The gateway then supplies a \term{first answer} with two components. The \term{event receipt} identifies the application and model versions active for the event; any output, score, decision, notice, warning, transmission, use, or intervention recorded for that event; and whether a requested function did not run. The \term{custodian map} identifies entities that supplied, configured, integrated, or operated material components; the holder and legal controller of each listed record; and the roles with authority to update, restrict, roll back, suspend, or stop a relevant function. The operator preserves records within its control, identifies the retention schedule and known deletion, and activates each contractual retrieval and preservation path required before deployment.

This answer gives the affected person identity and provenance without requiring reconstruction of vendor architecture. The operator forwards the notice and common event identifier to covered control holders so they can preserve and acknowledge their lanes. Local provenance and a retrieval right remain conditions of covered deployment even when a supplier relationship is confidential. Forwarding creates neither statutory coverage, choice of law, personal jurisdiction, nor service of process; direct duties attach only within the Act's valid reach, and other suppliers respond through the operator's contract or another independent source of law.

Stage Two begins only after that record exists and only by agreement or in a proceeding with an independent basis to order production. It is not a freestanding bill of discovery. The affected person must identify a claim-specific lane and proportionate information tied to the event and a merits issue; the next Section determines which covered enterprise, if any, remains.

\subsection{Material Control Requires Risk, Authority, and Information Materiality}
\label{sec:material-control}

``Ability to affect the system'' is too broad. A cloud provider, investor, data vendor, researcher, model author, and downstream institution may each have some causal connection to a deployment. Beyond Stage One, the rule should retain a covered control holder only when three requirements are satisfied at the relevant time.

First, \term{risk specificity}: the asserted control must connect to the risk alleged in the verified notice. General influence over corporate strategy, compute availability, or model research does not establish control over age access, hiring rank, clinical escalation, a particular tool permission, or another specified hazard.

Second, \term{practical authority}: the enterprise must have possessed an operational or legally enforceable ability to configure, constrain, evaluate, monitor, update, suspend, roll back, or terminate the relevant function. A theoretical technical capacity, an expired contract right, or influence that required another actor's voluntary cooperation is evidence but not practical authority by itself. Contractual audit, update, safety-override, and suspension rights matter because they may be usable controls; a contractual label declaring the parties independent does not decide the question.

Third, \term{information materiality}: the Stage One record must support that the function formed part of the event path or governed a safeguard, and the requested information must bear non-de-minimally on an element, defense, or remedy. An incidental contribution or a remote ability affecting every customer in the same undifferentiated way falls outside. The inquiry asks whether the requested information can materially inform adjudication. It does not require a showing that exercising the authority would have prevented the injury, changed the ultimate decision, or satisfied causation.

The requirements are cumulative. Nonexclusive factors include who selected the deployment's purpose and population; chose the model and version; configured instructions, memory, retrieval, data fields, tools, and credentials; set engagement, ranking, decision, or escalation thresholds; received evaluations, complaints, and incident signals; could issue an update or impose a restriction; retained contractual audit or safety-override rights; and could suspend, roll back, or terminate the function. The inquiry is temporal. Authority obtained after the event does not create earlier control, although later records may reveal what authority existed.

Several negative rules keep the test bounded. Investment, ordinary cloud capacity, publication of research, past employment, trademark association, and model provenance alone do not establish material control. Publishing open weights may end the original provider's ability to monitor or stop an independent deployment. A downstream modifier that removes safeguards, adds credentials, or selects a new high-impact purpose can occupy the relevant lane instead. Conversely, a provider that retains mandatory updates, remote suspension, risk monitoring, or binding deployment conditions may retain a lane despite contract language assigning ``all responsibility'' downstream.

Material control determines information obligations, not substantive liability. An enterprise can possess a control relevant to feasibility yet owe no duty under governing law, or receive an alert that does not constitute legal notice. The \term{merits firewall} therefore bars an inference based solely on coverage, notice, compliance, noncompliance, or participation. Independently relevant operational conduct retains only the effect governing law gives it.

\subsubsection{Adjudicating Material Control}

The test becomes useful doctrine when each requirement has an evidentiary question and an allocated burden. After the first answer, the affected person makes a prima facie showing by identifying a concrete lane connected to the asserted risk, showing that the Stage One record supports that the function formed part of the event path or governed a safeguard, and identifying requested information that bears non-de-minimally on a claim-specific element, defense, or remedy. The showing may point to an employer's enabled ranking feature, a provider's retained update authority, an application's self-harm threshold, or an integrator's workflow trigger. It concerns allocation and record relevance, not breach or causal sufficiency. An enterprise seeking termination then identifies Stage One material showing that the function did not run, no output was generated or transmitted, the asserted authority was unavailable at the relevant time, or the requested information bears no non-de-minimis relation to an element, defense, or remedy. The tribunal decides termination on that closed record or permits a bounded Stage Two inquiry when a genuine factual conflict remains.

The record duty supplies one limited burden consequence. If a covered enterprise's failure to create or retain a required Stage One field produces an ambiguity about its lane, the tribunal resolves that ambiguity against the record-duty holder only for access and termination. The consequence opens bounded production; it creates no inference of duty, breach, causation, discriminatory operation, or damages. That allocation prevents record noncompliance from becoming a safe exit while preserving the merits firewall.

\begin{table}[htbp]
\centering
\caption{Material Control as an Adjudicable Test}
\label{tab:material-control-proof}
\small
\begin{tabularx}{\textwidth}{>{\raggedright\arraybackslash}p{0.18\textwidth} >{\raggedright\arraybackslash}p{0.35\textwidth} X}
\toprule
Requirement & Evidence that supports the lane & Evidence that supports termination of the lane \\
\midrule
Risk specificity & Function map; risk ownership; feature requirements; evaluation scope; incident category; warnings; escalation assignment. & General corporate influence; unrelated research; infrastructure service without risk-specific configuration; records addressing a different feature or hazard. \\
Practical authority & Contract right; administrative permission; release or update role; audit access; monitoring feed; restriction, rollback, suspension, or termination mechanism; consistent exercise in practice. & Expired or legally unavailable right; access requiring another actor's uncommitted consent; technical ability prohibited and unused; independent modification beyond visibility or reach. \\
Information materiality & Stage One evidence that the function formed part of the event path or governed a safeguard; requested information bearing non-de-minimally on an element, defense, or remedy. & Feature inactive for the event; output never generated or transmitted; lane acquired after the event; requested information unrelated to the asserted path or claim-specific issue. \\
\bottomrule
\end{tabularx}
\end{table}

Reserved rights count when they were legally available and operationally usable, not merely written on paper. A right to obtain validation, require an update, audit use, or suspend a feature may qualify even if unused; a generic right to terminate an entire agreement may be too remote. Conversely, actual practice---mandatory updates, cross-tenant disabling, monitoring, or customer control over whether a score is used---may reveal authority that a contract label obscures. The applications in Part~III test those distinctions across integrated and divided deployments.

A material-control ruling terminates only the special front-door role. An entity never named on an underlying claim returns to ordinary nonparty status after completing any first-answer duty. A party already facing an independently sufficient cause of action remains subject to that claim even if it lacks the control lane alleged in the notice. This distinction gives safe exit its intended economy without using an information proceeding to decide a merits question it was not designed to reach.

\subsection{Creation, Preservation, and Production Are Separate Duties}
\label{sec:three-duties}

The front door depends on three legal functions that begin at different times.

\subsubsection{Record Creation Before Operation}

The \term{deployment-record duty} arises before a covered deployment affects the public. The visible operator has a nondelegable local duty to identify the system, version, purpose, population, material components, vendors, data sources, tools, credentials, intended limits, risk owners, release approval, event identifier, and intervention authority. ``Nondelegable'' here governs only compliance with the Act's record, retrieval, and access duties; it creates no underlying duty of care, vicarious liability, or primary-recourse rule. During operation, the operator records material changes, evaluations, complaints, incidents, escalations, overrides, restrictions, rollbacks, and termination decisions. A covered control holder maintains the subset corresponding to its lane. A supplier outside the Act's direct reach responds through the predeployment contract or another independent source of law.

A divided deployment requires a retrieval architecture as well as a list of vendors. Before use, the visible operator must obtain contractual rights to retrieve the identity, version, configuration, event, and control-allocation fields assigned to an upstream holder on demand. The agreement identifies a responsive agent, stable deployment and event identifiers, preservation and response duties, lawful routes for producing customer-controlled data, successor obligations after acquisition, and the treatment of subcontractors. The operator need not warehouse a supplier's source code or unrelated customer data. It keeps a closed local record sufficient to identify what ran and a legally usable path to obtain the defined supplier fields.

That requirement builds a usable retrieval right into deployment design. Where Rule 34 control turns on a legal right to obtain documents on demand, a covered operator can create that right before deployment instead of litigating it after harm. Where privacy law or a customer agreement places legal control with the operating institution, the contract can authorize a coordinated response rather than allowing the host and customer to direct the requester to each other. Operator notice alone does not create personal jurisdiction over a foreign or out-of-state supplier.\lrfootnote{See \casecite[264--68]{bristolMyers2017}; \casecite[359--71]{fordMotor2021}. The model statute therefore distinguishes a control holder directly subject to the Act from a supplier reached through the operator's contract and local record.} A domestic operator preserves the minimum record locally and conditions covered use on a reachable agent and retrieval duty.

The designation supplies the detail. A statute need not require indefinite retention of every prompt or every user's data. It should define event categories, sampling where appropriate, privacy minimization, and a retention period calibrated to the protected interest and ordinary limitations period. High-impact decisions affecting employment or benefits may require the inputs, output, version, governing policy, notice, and review event for each affected person while permitting aggregation or deletion of unrelated interaction content. Relational deployments involving minors may require crisis detections and interventions while minimizing intimate content unrelated to the asserted event.

\subsubsection{Preservation After a Defined Trigger}

Retention and litigation preservation are related but distinct. The statute fixes a baseline period before any dispute. A verified front-door notice, regulator demand, filed claim, or other event making a dispute reasonably foreseeable suspends ordinary deletion for the risk-specific record. The operator must issue a hold to relevant custodians and notify identified control holders. The notice should not freeze every model and customer record. Its scope follows the event, versions, period, user or decision, control lanes, and asserted risk.

This statutory trigger creates a stable bridge to existing preservation law. It also protects defendants: a dated notice and documented hold define what had to be preserved and when. The record can show that a missing item was outside scope, deleted before any duty attached, or available from another source.

\subsubsection{Staged Production After the Threshold}

Production proceeds in stages. Stage One is the first answer described above: identity and active versions; event-specific run, output, notice, and intervention fields; suppliers, custodians, and legal controllers; listed operational and stopping authority; and preservation, deletion, and retrieval status. Its purpose is to identify the deployment and test the control lanes.

Stage Two is not an independent production action. By agreement or in a pending proceeding with authority over the record, it follows only when the affected person uses Stage One to make a prima facie showing of a claim-specific control lane, information materiality, and proportional need. It can include evaluations and validation; incident trends; configuration history; approvals; rejected or conditional mitigations; escalation and response; and records needed to test an alternative or defense. Time, version, population, feature, and risk bound the request.

Source code, model weights, unrestricted training data, security-sensitive controls, and unrelated users' content occupy a final tier. A court should order them only when the requesting party shows that a material issue cannot reasonably be tested through less intrusive evidence such as versioned documentation, validation results, representative testing, counterfactual interrogation, interrogatories, an expert protocol, or in camera review.\lrfootnote{See \bluecite[1076--79]{stewartBeyondExplanation2026}. Counterfactual interrogation tests the behavior of a known decision-time version; Stage One supplies the identity and preservation predicates needed to direct that tool.} The distinction protects legitimate secrecy without converting a trade secret into a barrier to the legal process that governs its consequential use.

\subsection{Protection, Remedies, and Safe Exit}
\label{sec:remedies}

Confidentiality is part of both stages. Every disclosure remains limited to nonprivileged matter and may use redaction, pseudonymization, data minimization, limited access, secure expert review, in camera review, and protective orders.\lrfootnote{See \statcite{frcp26cTradeSecret}.} A protected record must still be identified by custodian, date, version, and asserted basis for withholding; a privilege log separates lawful protection from an unreviewable proprietary claim.

The statute should specify graduated consequences.

\begin{enumerate}[leftmargin=2.2em]
\item \term{Late or incomplete identification}: an order enforcing the statutory first answer and reasonable expenses caused by noncompliance. For a claim created or governed by the same enactment, the limitations period is suspended from service of a compliant request until a fixed period after compliance or final denial. Other state and federal claims retain their own limitations and tolling law.
\item \term{Failure to create or retain the minimum record}: the administrative penalty, corrective record order, compliance monitoring, or statutory-access remedy specified by the sectoral enactment. The consequence addresses the independent deployment obligation; it does not become a prescribed inference on the underlying injury.
\item \term{Loss after a preservation trigger}: an order to identify the loss and available substitute sources, followed by measures under the governing forum's preservation law. For electronically stored information that should have been preserved in anticipation or conduct of federal litigation and cannot be restored or replaced, Rule 37(e) supplies the federal standard. The state statute should not prescribe a competing litigation sanction for the same loss.
\item \term{Intentional concealment or falsification}: specified civil penalties, enforcement expenses, and remedies for the independent access violation. Litigation misconduct remains subject to the tribunal's governing authority.
\end{enumerate}

Safe exit is the reciprocal limit. After Stage One, an enterprise may move to terminate its front-door role by showing that it lacked material control over the asserted risk during the relevant period. The claimant may use the first-answer record to contest that showing. An entity not otherwise sued returns to ordinary nonparty status. A party remains only when an independently sufficient claim, a valid subpoena, or regulatory authority supplies a basis for continued participation. The gateway obligation therefore ends at the point where the control record shows that another lane or no AI-specific lane is relevant.

\subsection{Model Statutory Text}
\label{sec:model-text}

The following text consolidates the rule as a default model. A sectoral enactment would identify the covered field, retention clock, administering authority, and ordinary civil-penalty ceiling while preserving the sequence and merits firewall.

\begin{quote}
\small
\textbf{Section 1. Designated Deployment and Covered Control Holders.} This Act applies only to an AI deployment category designated by the enacting legislature or by a sectoral agency acting under express and bounded statutory criteria. The designation shall identify the covered function, protected interests, visible operator, affected persons, minimum record, retention period, material-change threshold, administering authority, enforcement route, and civil-penalty ceiling. A \emph{covered control holder} is a public or private organization to which this State's law validly applies that operates, provides, or continuously supports the identified covered deployment, or that has contractually assumed the response duties defined here. Before designation, the legislature or authorized agency shall publish a federal-overlap determination identifying applicable federal actors, statutes, regulations, exclusive enforcement provisions, express or conflict-preemption issues, excluded actors and fields, the relation between remedies, and the claims to which state tolling may lawfully apply.

\textbf{Section 2. Deployment Record and Retrieval Architecture.} Before making a covered deployment available or incorporating it into an affected process, the visible operator shall create and maintain a local record identifying the system and versions; purpose and affected population; material model, data, retrieval, tool, and credential components; deployment and event identifiers; suppliers and integrators; risk and release owners; evaluation and approval; material limitations and unresolved risk decisions; incident and escalation channels; and authority to update, restrict, roll back, suspend, or terminate. Each covered control holder shall maintain the corresponding portion and provide current provenance and a response contact. The operator shall obtain a contractual right to retrieve the defined supplier fields, require preservation and response, and carry those duties through subcontracting, acquisition, and service termination. These operator duties are nondelegable solely as to compliance with this Act's record, retrieval, and access requirements; they create no underlying nondelegable duty of care, vicarious liability, or primary-recourse rule. The designation supplies the retention period and material-change rule.

\textbf{Section 3. Stage One Notice; No Condition Precedent.} An affected person may deliver to the visible operator's designated agent a sworn declaration or declaration under penalty of perjury identifying: (a) an adverse event affecting an interest protected by the law named in the designation; (b) the covered deployment or visible operator through which the event occurred; (c) facts reasonably connecting the requested identity and event fields to that occurrence; (d) the event, transaction, account, application, decision, or period with as much specificity as reasonably available; (e) the closed fields sought; and (f) a good-faith intended use of the answer to evaluate, contest, appeal, or seek relief for the event. Stage One does not require a fact held exclusively in the deployment record. Counsel, if any, shall certify to the operator proper purpose and reasonable inquiry. Delivery invokes the first-answer duty but is not service of process for an enforcement action. Invocation or exhaustion of Stage One is not a condition precedent to commencing or maintaining any action other than an action to enforce the duties created by this Act. Failure to invoke Stage One, or an incomplete first answer, creates no pleading requirement, ground for dismissal, stay, or defense as to any other claim.

\textbf{Section 4. First Answer and Preservation.} Within thirty days after delivery, the operator shall provide an event receipt and custodian map. The receipt identifies the application and model versions active for the event; any requested output, score, decision, notice, warning, transmission, use, or intervention; and whether a named function did not run. The map identifies each entity that supplied, configured, integrated, or operated a material component; the custodian and legal controller of each field; and each entity with update, restriction, rollback, suspension, or termination authority. The operator shall preserve its risk-specific record, identify known deletion or unavailability, activate its contractual retrieval and preservation rights, and promptly forward the notice and common event identifier to each identified covered control holder. Each such holder shall preserve and acknowledge its assigned lane and may respond within that lane. For any other supplier, the operator shall satisfy the retrieval obligation through Section 2's contract and local record. Operator notice alone creates neither statutory coverage, choice of law, service, personal jurisdiction, nor consent. One extension of fifteen days may be obtained for a stated technical, privacy, or security need. Section 7's protections apply to every disclosure under this Section.

\textbf{Section 5. Material Control.} A covered control holder remains for Stage Two only upon a prima facie showing that, at the relevant time: (a) its function was specifically connected to the asserted risk; (b) it possessed practical authority through configuration, constraint, evaluation, monitoring, update, restriction, rollback, suspension, termination, or a legally enforceable equivalent; and (c) the Stage One record supports that the function formed part of the event path or governed a safeguard, and the requested information bears non-de-minimally on an element, defense, or remedy. This Section requires no showing that exercising the authority would have prevented the injury, changed the decision, or satisfied causation. Investment, general cloud service, publication, research, brand association, provenance, or theoretical technical capacity alone is insufficient. Any other supplier is reached only through Section 2's contract, ordinary subpoena authority, or another independent source of law.

\textbf{Section 6. Objection, Burdens, and Enforcement.} The operator or a responding covered control holder may object by field and shall state facts supporting noncoverage, nonoperation, burden, privilege, privacy, security, or documented unavailability. A proprietary label or referral to an unnamed supplier is insufficient. A timely, specific objection stays disclosure of only the contested field pending prompt neutral review; every uncontested field remains due. An affected person may seek expedited enforcement in a court of competent jurisdiction or another lawful forum. The requester bears the burden on coverage, delivery, and the Stage One occurrence threshold. The objecting operator or covered control holder bears the burden on a field-specific objection based on facts within its control. After the first answer, the requester identifies a prima facie control lane; a covered control holder seeking termination identifies the Stage One record defeating risk specificity, practical authority, or information materiality. An ambiguity caused by that holder's violation of Section 2 is resolved against it only for access and termination and creates no merits inference.

\textbf{Section 7. Stage Two, Federal-Court Boundary, and Protection.} Further nonprivileged production may be obtained by agreement or ordered only in a pending action, arbitration, or administrative review with an independent basis to exercise authority over the producing person and record. Stage Two creates no freestanding bill of discovery. Production requires that the producing holder remain under Section 5 and that the requested production be necessary and proportionate. It shall be limited by time, version, affected population, feature, and asserted risk. Source code, weights, sensitive data, raw content, expressive deliberation, and security controls require written findings that less intrusive evidence cannot reasonably resolve a material issue.

The protections in this Section apply to every disclosure under Sections 4 and 7. The owner of contested protected material and any materially affected third party shall receive notice and an opportunity for prompt neutral review before disclosure. The tribunal shall protect privilege, privacy, trade secrets, research, association, expression, and security through minimization, redaction, pseudonymization, limited access, secure expert review, in camera review, or another lawful safeguard.

In an action pending in a court of the United States, the Federal Rules of Civil Procedure govern all Stage Two production and protection. After an action concerning the adverse event is pending in such a court, and except for relief on a first-answer violation completed before that action was commenced, no party may invoke Stage One to obtain fields from another party in that action. As to fields first sought from another party after the action is pending, the first-answer duty is satisfied through, and limited by, the Federal Rules of Civil Procedure. This Act creates no entitlement to production from a party beyond the scope of Rule 26(b)(1), advances no discovery before Rule 26(d), and enlarges no production or subpoena authority under Rules 34 or 45. The duties to create, retain, and preserve the record under Sections 2 and 4 apply without regard to whether an action is pending.

\textbf{Section 8. Merits Firewall and Termination.} Coverage, notice, compliance, noncompliance, and participation create no inference or presumption of duty, breach, defect, discriminatory operation, agency, causation, damages, or joint liability solely by virtue of this Act. Governing law supplies every element and defense and determines the independent relevance, if any, of operational conduct. A covered control holder that completes its applicable Section 4 duties and defeats a Section 5 requirement terminates its front-door role. An entity not otherwise subject to an independently sufficient claim resumes nonparty status. A valid subpoena, regulatory authority, or another source of law may require continued participation on its own terms.

\textbf{Section 9. Remedies and Limitations.} The lawful forum may issue an expedited compliance order and award reasonable actual expenses caused by an incomplete, late, or evasive answer. The administering authority may order creation or correction of the minimum record, compliance monitoring, and a civil penalty calibrated to duration, affected persons, record scope, and culpability up to the ceiling fixed in the sectoral designation. Knowing falsification, concealment, or repeated violation permits recovery of enforcement expenses and an enhanced penalty up to twice that ceiling. Bare nonperformance creates no private statutory damages; Section 8 governs its merits effect, and this Act creates no damages for the underlying adverse event. Measures for lost electronically stored information in federal litigation remain governed by Federal Rule of Civil Procedure 37(e). Only a claim created or expressly governed by this Act is tolled from delivery of a compliant request until sixty days after compliance or final denial.

\textbf{Section 10. Nonwaiver, Forum, Standing, and Federal Procedure.} The duties to create, retain, and provide the minimum record are nonwaivable as to affected persons, arise from a covered deployment, and exist whether a dispute is litigated, arbitrated, or never filed. Except as an element of a claim to enforce the first-answer duty itself, invocation or exhaustion of Stage One is not a condition precedent to commencing or maintaining any civil action or arbitration. Failure to invoke Stage One, or an incomplete first answer, creates no pleading requirement, ground for dismissal, stay, or defense as to any other claim.

A valid arbitration agreement may select the private enforcement forum under the Federal Arbitration Act; the first-answer clock continues independently, and completion of Stage One is not a condition precedent to commencing or continuing arbitration except as an element of a claim to enforce the first-answer duty itself. An agreement to which the administering authority is not a party does not restrict public enforcement brought in that authority's own name.

A claim to enforce a first-answer violation completed before a federal action was commenced seeks the compliance relief specified in Section 9. In federal court, the Federal Rules govern commencement, pleading, service, provisional relief, adjudication, and judgment. Section 7 governs any fields first sought from another party after a federal action concerning the adverse event is pending.

Nothing in this Act creates Article III standing, personal jurisdiction, or subject-matter jurisdiction; overrides immunity or privilege; binds a nonsignatory to arbitration; or alters federal rules governing pleading, service, joinder, discovery, amendment, protective orders, injunctions, and sanctions. Courts and arbitrators administer the duties created here; they do not derive them from ordinary discovery authority.

\textbf{Section 11. Territorial Reach, Federal Supremacy, and Severability.} An upstream enterprise has a direct duty only when it purposefully supplied, configured, operated, or continuously supported the identified deployment for covered use in this State, or another lawful basis for legislative authority exists, and application of this State's law satisfies territorial and choice-of-law rules. Coercive relief additionally requires valid service and an independent basis for personal jurisdiction. Appointment of the operator's response agent authorizes delivery under Section 3 only and creates no consent to general jurisdiction. This Act imposes no requirement to the extent federal law preempts it, requires no disclosure prohibited by federal privacy, security, communications, or confidentiality law, creates no claim solely to enforce a federal duty reserved to the United States, and does not regulate the United States absent congressional authorization and an applicable waiver. A tribunal shall sever a preempted actor, field, remedy, or application when the remaining designation can operate consistently with the legislature's stated purpose.
\end{quote}

The regulated acts are operating without the required record and refusing the resulting access right, not unsafe operation in the abstract. The model text makes that allocation enforceable while leaving the merits where governing law placed them.

\subsection{The Front Door Ends Where Federal Discovery Begins}
\label{sec:erie}

The front door's most durable contribution precedes any access dispute. No Federal Rule tells an enterprise what version, approval, evaluation, escalation, intervention, event, or retrieval record it must create while operating a covered deployment. Rule 34 permits requests for designated existing materials within a party's possession, custody, or control; it does not require the enterprise to create the record in the first place. Thus, even if a federal court declined to enforce a first-answer demand against another party after litigation began, the ex ante duty would still have performed the Article's indispensable work. The deployment would remain identifiable, its event and control record would exist, and ordinary federal discovery could reach that record subject to the Rules' own limits. \emph{Erie} therefore governs the access channel in a bounded set of federal cases. It does not determine whether a covered deployment must remain observable.

\emph{Berk v. Choy} makes the correct axis explicit: scope, not characterization. The Court reversed the dismissal of a diversity medical-negligence action that lacked the affidavit of merit required by Delaware law. An eight-Justice majority held that Rule 8 answered whether the action could be dismissed because the plaintiff had not supplied merits evidence at the outset; Justice Jackson concurred only in the judgment, locating the conflicts in Rules 3 and 12 instead. The opinions disagreed over the precise Federal Rule, but each treated the ``question in dispute'' as decisive. Once a valid Federal Rule answers that question, contrary state law gives way even if substantive under \emph{Erie}; at the Rules Enabling Act step, the displaced law's substantive purpose makes no difference. Stage One's substantive character is therefore insufficient by itself. Its stronger defense is noncollision: a fixed, transaction-linked operating and access duty asks a question the Federal Rules do not answer.\lrfootnote{See \emph{Berk v. Choy}, 607 U.S. 187, 191--200 (2026) (reversing dismissal of the diversity action and holding that Federal Rule of Civil Procedure 8 displaced Del. Code tit. 18, \S~6853); \emph{id.} at 200--12 (Jackson, J., concurring in the judgment) (locating the conflicts in Rules 3 and 12 rather than Rule 8); \emph{Hanna v. Plumer}, 380 U.S. 460, 469--74 (1965); \emph{Shady Grove Orthopedic Associates, P.A. v. Allstate Insurance Co.}, 559 U.S. 393, 398--410 (2010); \emph{Sibbach v. Wilson \& Co.}, 312 U.S. 1, 14--16 (1941).}

The distinction is structural. Delaware required a plaintiff to submit evidence to a court to commence or maintain a lawsuit. The front door adds nothing to a complaint, changes no dismissal standard, and creates no prerequisite to another claim. It asks what an operator must record during covered operation and what fixed receipt it owes after a compliant request concerning a covered event, whether the recipient later sues, arbitrates, appeals, seeks correction, or brings no proceeding. Model Act sections 3 and 10 make that line operative.

The discovery Rules answer later questions. Rules 26, 34, and 45 govern the timing, scope, and targets of party and nonparty discovery; Rule 37(e) governs covered loss of electronically stored information. Those Rules presuppose an action, a discovery relation, and an identifiable target or record. Rule 27, the closest pre-action candidate, authorizes perpetuation of identified testimony rather than investigation to determine whether a claim exists. None tells an enterprise what record to create during operation or what transaction-linked receipt it owes a person who may never litigate.\lrfootnote{See Fed. R. Civ. P. 26(b)(1), 26(d)(1), 27(a), 34(a), 37(e), 45(a)(2); \emph{Workman v. United States Postal Service}, 127 F.4th 237, 240--45 (10th Cir. 2025).}

Existing access laws illustrate the institutional form without resolving the \emph{Erie} question. The FCRA requires disclosure of a bounded consumer file; Florida requires pre-suit insurance information; California requires access to personnel records but suspends that route during related litigation. In \emph{American Bank v. City of Menasha}, the Seventh Circuit likewise held that a public-records request did not become discovery solely because the requester sought material for litigation, while warning that substance controls if a State disguises discovery as inspection. The front door follows that warning only if its fields are fixed before the dispute, its duty operates outside litigation, and claim-specific relevance and proportionality remain outside Stage One. The state pre-suit proceedings surveyed in Section~\ref{sec:ordinary-procedure} demonstrate administrability, but as court-authorized discovery they do not supply this proposal's federal-law form.\lrfootnote{See Fair Credit Reporting Act, 15 U.S.C. \S~1681g(a)(1)--(3) (2024); Fla. Stat. \S~627.4137(1) (2026); Cal. Lab. Code \S~1198.5(b)--(c), (n)--(o) (2026); \emph{American Bank v. City of Menasha}, 627 F.3d 261, 265--67 (7th Cir. 2010).}

\emph{Berk} itself supplies the affirmative model. In distinguishing \emph{Cohen v. Beneficial Industrial Loan Corp.}, the Court explained that there was no conflict because state law created a liability while Rule 23 governed disclosure to the court and notice to interested parties. The front door belongs on that side of the line when drafted as the model text provides. Coverage, a compliant notice, delivery, and nonperformance complete a state-law violation. An enforcement action seeks the prescribed event receipt and custodian map as relief on that claim; it does not ask the court to compel evidence for proving a different claim. Rule 8 governs the enforcement complaint, Rule 12 governs dismissal, Rule 65 and ordinary remedial rules govern interim relief, and the remaining Federal Rules govern adjudication. State law supplies the duty and remedy; federal law supplies the litigation machinery.\lrfootnote{See \emph{Berk}, 607 U.S. at 195--96 (distinguishing \emph{Cohen}); \emph{Cohen v. Beneficial Industrial Loan Corp.}, 337 U.S. 541, 555--56 (1949).}

Calling production ``relief'' is not sufficient by itself. Two objective features separate the front door from a free-standing bill of discovery. First, the designation fixes the Stage One universe before the event. Two affected persons presenting different liability theories receive the same identity, version, event, custodian, and control fields for the same kind of event, although either may request fewer fields. Second, the duty performs the same work when no merits suit follows. A field whose scope changes with a pleaded element, defense, relevance argument, or proportionality showing is not a Stage One field; it belongs to Stage Two or ordinary discovery. Fixed content and litigation-independent operation are therefore elements of the rule's boundary, not labels assigned during enforcement.

The model Act withdraws where a Federal Rule answers the question. Once an action concerning the event is pending in federal court, Stage One cannot become an additional route around Rules 26(d), 26(b)(1), or 34; fields first sought from another party proceed under the Rules. A completed prefiling refusal remains an accrued access claim if jurisdiction and Article III permit. Thus an applicant who requested and was denied the fixed receipt before filing may seek that statutory relief, while one who files first uses federal discovery against the record that model Act section 2 already required. Sections 7 and 10 state the boundary rather than leaving it to a savings clause.

The retrieval architecture likewise complements Rule 34. Model Act section 2 creates a contractual right to retrieve supplier fields but does not declare that right to be Rule 34 control. A legal-right circuit may treat the clause as supplying a right to obtain documents; a practical-ability circuit may consider the clause, response mechanism, and course of dealing. The federal court independently decides possession, custody, or control. State contract law creates the entitlement; Rule 34 determines its procedural consequence.\lrfootnote{Compare \emph{In re Citric Acid Litigation}, 191 F.3d 1090, 1107--08 (9th Cir. 1999) (applying a legal-right-to-obtain test), with \emph{Sergeeva v. Tripleton International Ltd.}, 834 F.3d 1194, 1200--02 (11th Cir. 2016) (affirming, in a proceeding under 28 U.S.C. \S~1782, a control finding supported by the affiliated entities' ordinary course of dealing and access to responsive information).}

The forum boundary follows. The principal adopting institution is a state legislature, and the ordinary enforcement forum is the one the sectoral designation supplies under state law. In compliant cases no tribunal is involved: the operator delivers the receipt, and any later federal merits action proceeds under ordinary discovery. A federal court encounters the state access claim only through diversity, supplemental, or another independent jurisdictional basis; Congress can adopt the same architecture for a federal cause of action, in which event \emph{Erie} does not arise. The state statute creates neither subject-matter jurisdiction nor personal jurisdiction. Article III independently requires a concrete injury. A federal plaintiff should allege that denial impaired a specified use---identification, correction, contest, appeal, or pursuit of relief---or caused reasonable actual expense. State courts apply their own standing law.\lrfootnote{See 28 U.S.C. \S\S~1331, 1332, 1367 (2024); \emph{TransUnion LLC v. Ramirez}, 594 U.S. 413, 425--31, 440--42 (2021); \emph{Spokeo, Inc. v. Robins}, 578 U.S. 330, 339--43 (2016); \emph{Havens Realty Corp. v. Coleman}, 455 U.S. 363, 373--74 (1982); \emph{Public Citizen v. United States Department of Justice}, 491 U.S. 440, 449--50 (1989); \emph{Federal Election Commission v. Akins}, 524 U.S. 11, 21--25 (1998); \emph{ASARCO Inc. v. Kadish}, 490 U.S. 605, 617--18 (1989).}

Record preservation and limitations follow the same allocation. State law may prescribe baseline creation and retention as conditions of covered operation and may extend the defined retention period after a qualifying notice. If electronically stored information is lost in anticipation or conduct of federal litigation, Rule 37(e) supplies the federal litigation measures; the state statute prescribes no competing adverse inference or case-management sanction. A State may toll a limitations period for a claim the Act creates or expressly governs while the operator answers. \emph{Walker v. Armco Steel Corp.} remains the relevant scope precedent: Rule 3 did not answer when state limitations law was satisfied, so the state rule applied. Rule 15 continues to govern amendment and relation back in a later federal action.\lrfootnote{See \emph{Walker v. Armco Steel Corp.}, 446 U.S. 740, 750--52 (1980); \emph{Berk}, 607 U.S. at 205 (Jackson, J., concurring in the judgment); Fed. R. Civ. P. 15(c), 37(e).}

This allocation does not depend on how lower courts extend \emph{Berk} to other state gatekeeping devices or on a federal court accepting parallel access against another party after litigation begins. The rule is narrower: a State may require a covered deployment to remain observable and a fixed receipt to issue outside federal party discovery; once that discovery governs the same request, the Federal Rules supply the exclusive channel. The front door creates the record and first object. It ends where federal discovery begins.
